\chapter{相干多频驱动与多维能量格点}
\label{app:multifrequency-chapter}

\label{app:multifrequency-coherent}

正文先把单一载频的物理链闭合，因为这时一个整数边带指标就足以看清“周期相位 $\to$ Bessel 振幅 $\to$ 能量边带”的机制。本附录只做线性叠加，不引入新的电子动力学近似。以下所有公式都从正文的一般特征线解~\eqref{eq:characteristic-exact-solution} 出发。

单色只是最简单的情形。若场具有一般相干频谱
\begin{equation}
 \bm E(\bm r,t)
 =\operatorname{Re}\int_0^\infty\frac{\dd\omega}{2\pi}
 \bm E_\omega(\bm r)\ee^{-\ii\omega t},
 \label{eq:general-coherent-spectrum}
\end{equation}
则线性特征相位对频率逐项相加。定义
\begin{equation}
 \Ftilde(\omega)
 =\int\dd s\,
 \hat{\bm v}_0\cdot\bm E_\omega(\bm r_{\rm e}(s))
 \ee^{-\ii\omega s/v_0},
 \qquad
 \beta(\omega)=-\frac{q_{\rm e}}{2\hbar\omega}\Ftilde(\omega),
\end{equation}
便有
\begin{equation}
 \boxed{
 \frac{f_{\rm out}}{f_{\rm in}}
 =\exp\!\left\{
 \int_0^\infty\frac{\dd\omega}{2\pi}
 \left[-\beta(\omega)\ee^{-\ii\omega t_c}
 +\beta^*(\omega)\ee^{+\ii\omega t_c}\right]
 \right\}.}
 \label{eq:multifrequency-eikonal}
\end{equation}
对于离散的 $N$ 个互相相干频率 $\omega_j$，上式因指数可加而严格分解为
\begin{equation}
 \frac{f_{\rm out}}{f_{\rm in}}
 =\prod_{j=1}^{N}
 \exp\!\left[-\beta_j\ee^{-\ii\omega_jt_c}
 +\beta_j^*\ee^{+\ii\omega_jt_c}\right].
 \label{eq:multicolor-product}
\end{equation}
对每个因子使用正文已经固定的“$n_j>0$ 表示电子从第 $j$ 个频率通道获得能量”约定~\eqref{eq:sideband-amplitude-gain-convention}，逐项得到
\begin{align}
\frac{f_{\rm out}}{f_{\rm in}}
&=\sum_{n_1=-\infty}^{+\infty}\cdots
\sum_{n_N=-\infty}^{+\infty}
C_{\bm n}
\exp\!\left[-\ii t_c\sum_{j=1}^Nn_j\omega_j\right],\\
C_{\bm n}
&=\prod_{j=1}^{N}
J_{n_j}(2|\beta_j|)
\ee^{\ii n_j\arg(-\beta_j)}.
\label{eq:multicolor-amplitude-product}
\end{align}
因此多重整数指标 $\bm n=(n_1,\ldots,n_N)$ 对应的能量变化为
\begin{equation}
\boxed{
 \Delta\mathcal E_{\bm n}
 =\sum_{j=1}^N n_j\hbar\omega_j.}
 \label{eq:multicolor-energy-lattice}
\end{equation}
这里不再存在额外的正负号约定：它与正文单色通道完全一致。若各 $\omega_j$ 彼此不共振，$\bm n$ 通常唯一标记一个能量格点；若存在整数关系，不同 $\bm n$ 可能落在同一总能量，此时这些路径必须在振幅层相干相加。这样双频 PINEM、光学拍频和频率梳驱动都只是同一特征线母式的线性推广。
